
\documentclass{article}
\usepackage{colortbl}
\usepackage{makecell}
\usepackage{multirow}
\usepackage{supertabular}

\begin{document}

\newcounter{utterance}

\centering \large Interaction Transcript for game `cladder', experiment `full\_v1.5\_default', episode 4597 with Qwen3.6{-}27B.
\vspace{24pt}

{ \footnotesize  \setcounter{utterance}{1}
\setlength{\tabcolsep}{0pt}
\begin{supertabular}{c@{$\;$}|p{.15\linewidth}@{}p{.15\linewidth}p{.15\linewidth}p{.15\linewidth}p{.15\linewidth}p{.15\linewidth}}
    \# & \multicolumn{2}{c}{Player} && \multicolumn{2}{c}{Game Master} \\
    \hline

    \theutterance \stepcounter{utterance}  
    & & & \multicolumn{4}{p{0.6\linewidth}}{
        \cellcolor[rgb]{0.9,0.9,0.9}{
            \makecell[{{p{\linewidth}}}]{
                \texttt{\tiny{[P1$\langle$GM]}}
                \texttt{You are an expert in causal inference. The following question is not a typical commonsense query, but rather a meticulously designed question created by a professor specializing in causal inference, intended to assess the students' mastery of the course content.} \\
\texttt{Imagine a self{-}contained, hypothetical world with only the following conditions, and without any unmentioned factors or causal relationships: Having a brother has a direct effect on wife and alarm clock. Wife has a direct effect on alarm clock. For people who do not have a brother and wives that don't set the alarm, the probability of ringing alarm is 88\%. For people who do not have a brother and wives that set the alarm, the probability of ringing alarm is 45\%. For people who have a brother and wives that don't set the alarm, the probability of ringing alarm is 58\%. For people who have a brother and wives that set the alarm, the probability of ringing alarm is 13\%. For people who do not have a brother, the probability of alarm set by wife is 9\%. For people who have a brother, the probability of alarm set by wife is 64\%. Does having a brother negatively affect alarm clock through wife?} \\
\texttt{Start your answer with 'Yes' or 'No', followed by additional reasoning or evidence to support your explanation.} \\
            }
        }
    }
     \\ \\

    \theutterance \stepcounter{utterance}  
    & \multicolumn{4}{p{0.6\linewidth}}{
        \cellcolor[rgb]{0.561,0.737,0.561}{
            \makecell[{{p{\linewidth}}}]{
                \texttt{\tiny{[P1$\rangle$GM]}}
                \texttt{Yes} \\
\\ 
\texttt{To determine whether having a brother negatively affects the alarm clock through the wife, we must calculate the **Indirect Effect** (also known as the Natural Indirect Effect) of "Having a Brother" (\$B\$) on "Alarm Ringing" (\$A\$) mediated by "Wife Setting Alarm" (\$W\$).} \\
\\ 
\texttt{The causal structure is:} \\
\texttt{\$B \textbackslash{}rightarrow W \textbackslash{}rightarrow A\$} \\
\texttt{\$B \textbackslash{}rightarrow A\$ (direct effect, which we ignore for the indirect path)} \\
\\ 
\texttt{We need to compare the expected probability of the alarm ringing if the wife's behavior were set to what it would be if the person had a brother, versus what it would be if the person did not have a brother, while holding the direct effect of \$B\$ on \$A\$ constant (or rather, isolating the path through \$W\$).} \\
\\ 
\texttt{The standard way to assess the sign of the indirect effect in this context is to look at the change in the mediator (\$W\$) caused by the exposure (\$B\$), and then how that change in \$W\$ affects the outcome (\$A\$).} \\
\\ 
\texttt{**Step 1: Determine the effect of \$B\$ on \$W\$ (Exposure \$\textbackslash{}rightarrow\$ Mediator)**} \\
\texttt{*   \$P(W=1 | B=0) = 0.09\$ (Probability wife sets alarm if no brother)} \\
\texttt{*   \$P(W=1 | B=1) = 0.6} \\
            }
        }
    }
    & & \\ \\

    \theutterance \stepcounter{utterance}  
    & & & \multicolumn{2}{p{0.3\linewidth}}{
        \cellcolor[rgb]{0.9,0.9,0.9}{
            \makecell[{{p{\linewidth}}}]{
                \texttt{\tiny{[GM$|$GM]}}
                \texttt{Yes} \\
            }
        }
    }
    & & \\ \\

    \theutterance \stepcounter{utterance}  
    & & & \multicolumn{2}{p{0.3\linewidth}}{
        \cellcolor[rgb]{0.9,0.9,0.9}{
            \makecell[{{p{\linewidth}}}]{
                \texttt{\tiny{[GM$|$GM]}}
                \texttt{yes} \\
            }
        }
    }
    & & \\ \\

    \theutterance \stepcounter{utterance}  
    & & & \multicolumn{2}{p{0.3\linewidth}}{
        \cellcolor[rgb]{0.9,0.9,0.9}{
            \makecell[{{p{\linewidth}}}]{
                \texttt{\tiny{[GM$|$GM]}}
                \texttt{game\_result = WIN} \\
            }
        }
    }
    & & \\ \\

\end{supertabular}
}

\end{document}
